% ===========================================================================
%  sviluppatore/04-motore/05-live.tex
% ===========================================================================
\section{Lo stesso calcolo, su un conto vero}\label{sez:live}
% Sezione fitta di tabelle che non si spezzano: con il fondo pagina
% allineato, lo spazio avanzato finiva attorno a una tabella. Qui il fondo
% resta libero; alla fine della sezione si torna all'allineamento del libro.
\raggedbottom

\pkgLive{} non è un secondo motore: è lo stesso, con un broker al posto dei
dati storici. Il valore di tutto il progetto dipende da questa frase, e il
modo in cui viene mantenuta vera merita di essere letto per intero, perché è
anche l'esempio migliore di che cosa va storto quando non lo è.

\subsection{Il difetto che ha definito il progetto}

Fino a un certo punto, il backtest normalizzava la volatilità di ogni stream
e l'operatività dal vivo no. Il commento in testa a
\file{live/ensemble\_live.py} dice che cosa significava:

\begin{quote}\itshape
Senza il punto~2 un libro dichiarato al \num{4}\,\% di vol veniva mostrato
in backtest moltiplicato per \num{2.5} ed eseguito dal vivo a dimensione
piena: \num{0.4} volte il rischio e il rendimento misurati. Al contrario, un
libro al \num{25}\,\% veniva tagliato a \num{0.4} in backtest e girava dal
vivo a \num{2.5} volte.
\end{quote}

\noindent
Nessun errore, nessuna eccezione, nessun avviso. Il backtest mostrava un
numero e il conto ne produceva un altro, e la differenza era un fattore fra
\num{0.4} e \num{2.5} a seconda della volatilità dello stream --- cioè più
grande quanto più il libro si discostava dall'obiettivo, che è esattamente il
caso in cui si ha bisogno che il backtest sia attendibile.

\begin{perche}{Perché il rimedio è stato condividere la funzione}
Si sarebbe potuto riscrivere la normalizzazione in \pkgLive{} copiando la
formula. Due implementazioni della stessa formula restano identiche finché
qualcuno non ne tocca una, e questo difetto è nato esattamente così. La
soluzione è che \simbolo{\_vol\_trailing} è \emph{una}, e le due strade la
usano in modo diverso in un solo punto documentato: il backtest la shifta di
un giorno, il live no (\vediSez{ensemble}).
\end{perche}

\subsection{I due punti in cui la normalizzazione entra}

\begin{enumerate}
  \item \textbf{I pesi si calcolano sui rendimenti scalati}, non su quelli
        grezzi. Le modalità di pesatura guardano la matrice di covarianza: su
        serie non normalizzate uno stream al \num{25}\,\% di volatilità
        riceve un peso del tutto diverso da quello che riceveva in backtest.
  \item \textbf{\simbolo{scales()} restituisce il fattore da applicare adesso}
        al vettore di posizioni --- cioè la leva che il backtest stava
        mostrando e che il live non applicava.
\end{enumerate}

\subsection{Il riscaldamento}

La volatilità trailing non esiste prima di una ventina di osservazioni, e fra
venti e sessanta è ancora rumorosa. In quel tratto il sistema mescola la
stima con una volatilità \emph{a priori} --- per uno stream di approccio, la
sua volatilità di backtest --- con peso lineare crescente.

\begin{center}
\begin{tabular}{@{}L{40mm}L{98mm}@{}}
\toprule
\textbf{Osservazioni} & \textbf{Che cosa si usa} \\
\midrule
\rigaTenue
meno di \cod{SCALE\_MIN\_OBS} & Solo la volatilità a priori. Senza una, il
  fattore è \cod{1.0}. \\
fra le due soglie & Miscela lineare fra priori e stima. \\
\rigaTenue
oltre \cod{scale\_min\_obs} & Solo la stima: qui coincide con il backtest. \\
\bottomrule
\end{tabular}
\end{center}

\noindent
Il requisito dichiarato è: \emph{mai un salto di dimensione in un solo
giorno}. Una transizione netta fra priori e stima produrrebbe, il giorno in
cui la soglia viene superata, un ordine di ribilanciamento che non
corrisponde a nessun cambiamento del mercato.

\begin{nota}{Lo stato sopravvive ai riavvii}
I rendimenti per stream sono salvati su JSON, così la finestra non riparte da
zero a ogni riavvio del processo. Senza, ogni riavvio rimetterebbe il sistema
in riscaldamento, e il riscaldamento usa la volatilità a priori: il conto
cambierebbe dimensione per il solo fatto di essere stato riavviato.
\end{nota}

\subsection{I broker}

L'interfaccia è \simbolo{live.broker.LiveBroker}, la stessa classe che
\pkgPorte{} ri-esporta come porta \simbolo{Broker}. Le implementazioni
attuali coprono Interactive Brokers (\file{live/brokers/ibkr.py}),
MetaTrader~5 (\file{live/brokers/mt5.py}), l'esecuzione su carta e un broker
di prova con il tempo accelerato --- quest'ultimo serve alla pagina di
simulazione, che fa passare una giornata di mercato in pochi secondi.

Oltre ai metodi obbligatori, cinque \emph{ganci} con un default innocuo:

\begin{center}
\begin{tabular}{@{}L{46mm}L{92mm}@{}}
\toprule
\textbf{Metodo} & \textbf{A che serve} \\
\midrule
\rigaTenue
\simbolo{account\_label()} & Una riga che dice a quale conto ci si è
  collegati: numero, demo o reale, terminale, cambio applicato. Va nel
  registro all'avvio. \\
\simbolo{account\_identity()} & Chi è il conto, in un dizionario
  (\cod{broker}, \cod{account}, \cod{where}, \cod{kind}); \cod{None} se il
  broker non sa dirlo. Lega la conferma degli ordini reali a un conto
  (\vediSez{live-ordini-reali}). \\
\rigaTenue
\simbolo{real\_orders\_blockers()} & I motivi per cui il conto
  \emph{non} può ricevere ordini da Lyra. Se la lista non è vuota, un
  runner con ordini reali non parte; se compare mentre gira, i cicli si
  saltano. \\
\simbolo{open\_orders()} & Gli ordini ancora in lavorazione,
  \cod{\{simbolo: quanti\}}. Il motore non manda un ordine su un simbolo che
  ne ha uno in corso. \\
\rigaTenue
\simbolo{cancel\_open\_orders()} & Cancella gli ordini aperti \emph{di
  Lyra} (non quelli fatti a mano) e dice quanti. Il runner la chiama
  all'avvio e all'arresto. \\
\bottomrule
\end{tabular}
\end{center}

\noindent
L'esito di ogni ordine è un \simbolo{Fill} con \cod{status} fra
\cod{filled}, \cod{partial}, \cod{pending} (accettato, né eseguito né
cancellato), \cod{skipped} (sotto la quantità minima, non inviato),
\cod{rejected}, \cod{cancelled} ed \cod{error}, e una \cod{note} con il
perché nelle parole del broker. Un adattatore che non riesce a leggere il
conto solleva \simbolo{BrokerUnavailable} (una \simbolo{ConnectionError}):
il runner lo riprova invece di fermarsi. Anche un fill parziale
entra nel libro della strategia (\simbolo{LivePortfolioMirror.record\_fill}).

\begin{center}
\begin{tabular}{@{}L{30mm}L{108mm}@{}}
\toprule
\textbf{Adattatore} & \textbf{Che cosa garantisce} \\
\midrule
\rigaTenue
IBKR & L'ordine va al conto indicato (\cod{order.account}) e porta il
  riferimento del runner (\cod{orderRef} \cod{lyra:<id>}); equity e
  posizioni si leggono di quel conto e solo per i contratti \cod{STK}, così
  un'opzione su AAPL non viene letta come azioni AAPL. Un ordine non
  eseguito entro \cod{fill\_timeout\_s} viene cancellato; se la
  cancellazione non arriva entro \cod{cancel\_timeout\_s} l'esito è
  \cod{pending}, e un ordine \cod{Inactive} viene cancellato perché può
  riattivarsi. \simbolo{equity\_in\_valuta()} converte la
  \cod{NetLiquidation} della valuta base del conto con il tag
  \cod{ExchangeRate} (in riserva il cambio dei dati di Lyra): senza, un
  conto in euro risultava a equity zero. \simbolo{open\_orders()} legge
  \simbolo{reqAllOpenOrders()} (solo \cod{STK}, solo il conto);
  \simbolo{cancel\_open\_orders()} cancella solo i riferimenti
  \cod{lyra}. Il collegamento passa \cod{account} e un \cod{timeout} di
  \num{20}\,s. Posizioni, valori del conto e stati degli ordini di
  \pkg{ib\_insync} si aggiornano solo mentre il suo event loop gira, e fra
  un ciclo e l'altro il runner dorme: prima di leggere il conto
  \simbolo{\_aggiorna()} fa girare il loop e le posizioni si chiedono al
  server (\simbolo{reqPositions()}), altrimenti un ordine eseguito
  durante la notte risultava ancora da fare; se la richiesta fallisce
  l'errore sale, senza ripiegare sulla copia locale. Le richieste bloccanti
  hanno \cod{RequestTimeout} di \num{20}\,s (di default aspettano per
  sempre). Si importa \pkg{ib\_async}, il seguito mantenuto, oppure
  \pkg{ib\_insync}, nel thread che lo usa. \\
MT5 & Il riempimento segue \cod{symbol\_info.filling\_mode} (IOC, FOK o
  Return): imporre IOC fa rifiutare ogni ordine sui simboli che non lo
  accettano. Il retcode \num{10010} è un'esecuzione parziale, il
  \num{10008} un ordine accettato non ancora eseguito (\cod{pending}). Un
  conto \emph{hedging}, il trading disabilitato, \cod{trade\_expert}
  falso, \voce{Algo Trading} spento o \cod{tradeapi\_disabled} bloccano gli
  ordini reali. Un conto in un'altra valuta si converte con il tick della
  coppia nel terminale (\cod{EURUSD}, anche con il suffisso del broker).
  Gli ordini in attesa vengono da \simbolo{orders\_get()}; si cancellano
  con \cod{TRADE\_ACTION\_REMOVE} solo quelli con il \cod{magic} di Lyra.
  \cod{None} da \simbolo{positions\_get()}, \simbolo{orders\_get()} o
  \simbolo{account\_info()} è un errore dell'API, non un conto vuoto: letto
  come «nessuna posizione» faceva ricomprare il libro.
  \simbolo{is\_connected()} interroga il terminale, così dopo un suo
  riavvio ci si ricollega. Il pacchetto parla con un terminale per
  processo: le chiamate passano da un lucchetto, e \simbolo{shutdown()} lo
  chiama solo l'ultimo che si scollega. \\
\bottomrule
\end{tabular}
\end{center}

\begin{attenzione}{Le credenziali restano su file locale}
\percorsoHTTP{/api/brokers/credentials} legge e scrive un file sul computer
dell'utente. Non passano dal sito, non finiscono nel database, e non c'è
nessun percorso che le porti in rete. Questo è il motivo per cui la porta
della dashboard non va esposta: chi la raggiunge può leggerle. Per IBKR
porta e client ID sono controllati al salvataggio, e
\simbolo{live\_runners.cfg\_ibkr()} li passa al collegamento come
\cod{broker\_cfg}.
\end{attenzione}

\subsection{Simulazione o ordini reali}\label{sez:live-ordini-reali}

Un runner su un broker vero parte in simulazione: \cod{dry\_run} vero, il
motore calcola gli ordini e non li invia. Gli ordini reali richiedono tre
cose insieme, e nessuna basta da sola:

\begin{codiceInterno}{python}{Quando un runner invia davvero}{dashboard/live\_runners.py:216}
def ordini_reali(runner: dict) -> bool:
    r = runner or {}
    return (str(r.get("broker") or "") in BROKER_REALI
            and r.get("dry_run") is False
            and bool(r.get("real_orders_confirmed_at")))
\end{codiceInterno}

\noindent
\cod{dry\_run} deve essere \cod{False} in modo esplicito --- non basta che
manchi --- e la conferma la scrive solo una richiesta che la porta:
\cod{confirm\_real\_orders} quando si salva il runner, \cod{confirm} quando
si cambia modalità. L'avvio di un runner con ordini reali chiede a sua volta
\cod{confirm\_real}. Le rotte \percorsoHTTP{/api/live/*}, rimaste per le
prove da script, su un broker vero girano sempre in simulazione: gli ordini
reali passano solo dai runner, dove ci sono conferma, registro e kill-switch
salvato.

\subsubsection{La conferma è legata a un conto}

Al primo avvio dopo una conferma, \simbolo{\_prepara\_conto} scrive in
\file{data/runners/<id>/account.json} la chiave del conto
(\simbolo{chiave\_conto()}: broker, server per MT5, numero) insieme a
\cod{real\_orders\_confirmed\_at}. Agli avvii successivi con la stessa
conferma un conto diverso ferma l'avvio (\simbolo{\_NonParte}); a ogni ciclo
\simbolo{\_controlla\_conto} confronta di nuovo l'identità e, se è cambiata,
ferma il runner prima del ciclo (\simbolo{\_Fermo}) senza toccare gli
ordini del conto nuovo. Una conferma nuova lega il runner al conto del
prossimo avvio (\simbolo{legame\_valido()}). Senza identità
(\cod{account\_identity()} a \cod{None}), o con \file{account.json}
illeggibile (\simbolo{leggi\_legame()}), gli ordini reali non partono.

\subsection{Il registro del runner}\label{sez:live-registro}

Ogni riga scritta con \simbolo{REGISTRO.log(rid, testo, kind, data)} va in
memoria (le ultime \num{200}, per lo stato) e su disco, in
\file{data/runners/<id>/journal.jsonl}: un oggetto JSON per riga con
\cod{ts} (UTC), \cod{session} (una per avvio), \cod{mode}, \cod{broker},
\cod{kind}, \cod{msg} e, quando serve, \cod{data}. Oltre \num{5}\,MB il
file ruota su \file{journal.1.jsonl} \ldots{} \file{journal.9.jsonl}
(\simbolo{\_ruota\_registro}), e \simbolo{registro\_leggi} li ripercorre
all'indietro finché servono righe. I \cod{kind} sono \cod{mode},
\cod{start}, \cod{stop}, \cod{info}, \cod{cycle}, \cod{order} (simulati),
\cod{fill} (inviati, con l'esito), \cod{warning}, \cod{error}, \cod{risk} e
\cod{maintenance}.

\noindent
Nella stessa cartella stanno lo stato del motore (\file{live\_state.json}),
la storia della normalizzazione (\file{ensemble.json}), lo stato del
\simbolo{RiskManager} (\file{risk.json}) e il legame con il conto
(\file{account.json}). Prima stato e normalizzazione erano un file unico in
\file{data/}: due runner attivi insieme se li sovrascrivevano a vicenda.

\subsection{Dati vecchi}\label{sez:live-dati-vecchi}

\simbolo{LoaderLiveFeed} legge gli stessi file del backtest e non scarica
niente. Se nessuno li aggiorna, il motore continua a decidere sugli stessi
prezzi --- e con ordini veri dimensiona le quantità su prezzi di giorni
prima. \simbolo{live.freshness.motivo\_dati\_vecchi} conta le \emph{sedute
mancanti} fra l'ultima barra e l'ultimo giorno di borsa prima di oggi
(UTC), sul calendario del mercato del simbolo (\simbolo{calendario\_per},
dall'asset class). Prima di oggi, perché il feed legge i file con una data
di fine esclusiva: la barra di oggi la serve domani, e pretenderla dopo la
chiusura americana tratteneva ogni ordine dalle 22 a mezzanotte.

\begin{center}
\begin{tabular}{@{}L{30mm}L{78mm}L{28mm}@{}}
\toprule
\textbf{Calendario} & \textbf{Sedute} & \textbf{Tolleranza} \\
\midrule
\rigaTenue
\cod{nyse} & Lunedì--venerdì meno le festività NYSE (regole di
  \pkg{pandas}: Venerdì Santo, Juneteenth dal 2022, 4 luglio e Natale con
  il giorno osservato) & nessuna \\
\cod{fx} & Lunedì--venerdì & una seduta \\
\rigaTenue
\cod{crypto} & Tutti i giorni & una seduta \\
\bottomrule
\end{tabular}
\end{center}

\noindent
Mancano le chiusure straordinarie (lutti nazionali, uragani): quel giorno i
dati risultano vecchi di una seduta e gli ordini aspettano, cioè l'errore va
nella direzione prudente. Per barre infragiornaliere si contano le barre
mancanti in tempo di borsa (oltre tre, vecchi); per barre settimanali, oltre
due periodi e tre giorni.

\noindent
Il motore misura ogni simbolo a ogni ciclo, anche in simulazione
(\simbolo{\_dati\_vecchi}), e riporta il riassunto in
\cod{CycleReport.data\_stale}. Fuori dal dry-run trattiene gli ordini dei
simboli indietro e lo scrive in \cod{orders\_held}: un'ultima barra fresca
su SPY non rende freschi i prezzi di AAPL. Un simbolo del runner con una
posizione aperta e nessun prezzo nel ciclo (file mancante, simbolo saltato
dal caricamento) finisce nello stesso riassunto: il motore non può né
valutarlo né ridurlo. In entrambi i casi gli ordini che aumentano
l'esposizione aspettano (vedi \emph{Prima le riduzioni}, più avanti). Passano solo le chiusure
ordinate dal kill-switch del conto, che vendono la quantità che c'è e non
dipendono dal prezzo. Il controllo vale per ogni broker tranne quello su
carta (\cod{guard\_stale\_data}), che il simulatore TEST e la prova da
terminale guidano con il proprio orologio.

\subsection{Il kill-switch del conto}\label{sez:live-kill}

Il \simbolo{RiskManager} vive quanto il processo, e il livello del conto
protegge del denaro che sta nel conto. Per questo, con ordini reali, picco
di equity e chiusura sono salvati per \emph{conto}, in
\file{data/accounts/<chiave>.json} (\simbolo{carica\_stato\_conto},
\simbolo{salva\_stato\_conto}); mai nella cartella temporanea, che il
sistema svuota. Il resto dello stato --- strategie sospese con la data,
asset class ridotte con la scadenza, finestre di PnL --- va in
\file{risk.json} del runner con il suo contesto (modalità, broker, conto),
attraverso \simbolo{export\_state()} e \simbolo{restore\_state()}: uno
stato per contesto (\simbolo{contesti\_rischio}), così una corsa in
simulazione non cancella le strategie sospese della modalità reale.

\begin{itemize}
  \item Il drawdown si controlla all'inizio del ciclo, appena letta
        l'equity (\simbolo{check\_account}), non solo dopo gli ordini: un
        calo avvenuto fra due cicli faceva partire un altro giro di acquisti
        prima che il kill-switch scattasse. Con il conto chiuso il ciclo
        procede anche con l'equity a zero: chiudere non ne ha bisogno.
  \item Passare a simulazione e tornare agli ordini reali, o usare un altro
        runner sullo stesso conto, non riapre il conto: le rotte di
        creazione e di cambio modalità non toccano più lo stato. Lo riapre
        solo \rotta{POST}{/api/runners/<rid>/kill-switch/reset}, che azzera
        il file del conto a cui il runner è legato.
  \item Uno stato del conto illeggibile ferma gli ordini reali: trattarlo da
        vuoto riaprirebbe un conto chiuso.
  \item Se all'avvio l'equity è già sotto il picco del conto oltre la
        soglia, il kill-switch scatta e si scrive nel file del conto, ma il
        runner non manda ordini: non può sapere se è un prelievo o una
        perdita. Al prossimo avvio chiude (perdita); il reset lo evita
        (prelievo). Con l'equity non leggibile all'avvio non parte.
  \item Un conto chiuso non impedisce l'avvio: il runner parte in sola
        chiusura (il \simbolo{RiskManager} riparte chiuso, ogni peso è
        zero), chiude e a conto piatto (\simbolo{\_posizioni\_aperte}) si
        ferma; in sola chiusura nessun errore lo ferma. L'elenco dei runner
        segnala la sola chiusura (\cod{close\_only}) solo se il conto chiuso
        è quello della conferma attuale; altrimenti \cod{killed\_unbound}.
  \item Quando il kill scatta durante un ciclo, il ciclo di chiusura parte
        subito, senza aspettare la cadenza; se resta qualcosa di aperto si
        riprova dopo \cod{min(passo, 60\,s)} per tre volte, poi alla
        cadenza normale.
\end{itemize}

\noindent
Il drawdown del conto si controlla a ogni ciclo
(\simbolo{check\_account}); i livelli per strategia e per asset class, che
lavorano su finestre di barre, una volta per barra nuova
(\simbolo{\_nuova\_barra\_rischio}), come nel backtest: aggiornarli a ogni
ciclo --- quattro per barra --- accorciava la finestra di \num{30} giorni a
\num{35} cicli. Il PnL realizzato fra due barre si attribuisce tutto alla
barra nuova (\simbolo{LivePortfolioMirror.consolida\_realizzati}). Il
livello per asset class riceve dal ciclo il PnL di prezzo per classe una
volta per barra (\simbolo{\_registra\_pnl\_asset\_class}). Nel backtest
\simbolo{update} chiama \simbolo{check\_account} con la stessa regola di
prima: i risultati non cambiano.

\subsection{Il ciclo dal vivo}

Il runner fa il primo ciclo subito, poi uno ogni quarto della barra più
corta che guarda (\simbolo{cadenza\_reale\_secondi}, mai sotto i \num{15}
secondi): per strategie giornaliere, ogni sei ore. L'avvio dalla pagina
imponeva \num{60} secondi, e un runner giornaliero rifaceva lo stesso ciclo
ogni minuto.

\noindent
Ogni runner gira in un thread suo, con un proprio event loop
\cod{asyncio}: \pkg{ib\_insync} ne ha bisogno per collegarsi, e un thread
che non è il principale non ne ha uno. Fermato il runner, gli ordini di Lyra
ancora aperti si cancellano (con ordini reali, sul conto verificato), il
collegamento al broker si chiude --- altrimenti IBKR terrebbe occupato il
client ID --- e il loop con lui.

\begin{itemize}
  \item \textbf{Un avvio solo.} \simbolo{start} tiene un lucchetto dai
        controlli fino alla registrazione del runner: con due richieste
        ravvicinate potevano partire due motori sullo stesso conto. Un
        errore dopo la registrazione rimette il runner fermo.
  \item \textbf{Equity assente.} Un conto con equity zero o mancante salta
        il ciclo (\cod{skipped\_reason} \cod{no\_account\_equity}): prima
        ogni peso andava a zero e partiva una chiusura per ogni posizione.
  \item \textbf{Ordini in corso.} Un simbolo con un ordine ancora aperto al
        broker non riceve un ordine nuovo; se gli ordini aperti non si
        leggono, non parte nulla.
  \item \textbf{Prima le riduzioni.} Gli ordini che riducono una posizione
        (\simbolo{\_riduce}) partono prima di quelli che la aumentano, e
        un'inversione diventa chiusura più apertura
        (\simbolo{\_separa\_inversioni}; la gamba di apertura conta come
        aumento). Se qualcosa è trattenuto (dati vecchi, ordine in corso,
        posizione senza prezzo) o una riduzione non è \cod{filled} né
        \cod{skipped}, gli aumenti del ciclo non partono: in una rotazione
        con la vendita trattenuta, l'acquisto da solo portava l'esposizione
        oltre quella del libro vecchio e di quello nuovo.
  \item \textbf{Stop fra un ordine e l'altro.} Il runner assegna al motore
        \cod{stop\_requested}: premuto Stop, gli ordini non ancora inviati
        del ciclo in corso non partono.
  \item \textbf{Cicli falliti.} Un errore di collegamento
        (\simbolo{e\_transitorio}: \simbolo{ConnectionError} e quindi
        \simbolo{BrokerUnavailable}, \simbolo{TimeoutError},
        \simbolo{OSError}) si riprova senza limite, con attese che
        raddoppiano fino a \num{15} minuti. Un'eccezione di altro tipo si
        riprova dopo \cod{min(passo, 60\,s)}; tre di fila fermano il runner
        con la causa, tranne in sola chiusura. Un ciclo saltato non tocca né
        il kill-switch né la manutenzione.
\end{itemize}

\begin{center}
\begin{tabular}{@{}L{50mm}L{88mm}@{}}
\toprule
\textbf{Rotta} & \textbf{Che cosa fa} \\
\midrule
\rigaTenue
\rotta{GET}{/api/runners} & I runner salvati, con \cod{mode},
  \cod{real\_orders}, \cod{killed}, \cod{close\_only}, il conto legato,
  \cod{data\_stale}, \cod{account\_blocked} e le soglie del
  kill-switch. \\
\rotta{POST}{/api/runners} & Crea o aggiorna. Ordini reali solo con
  \cod{dry\_run: false} e \cod{confirm\_real\_orders: true}. \\
\rigaTenue
\rotta{POST}{/api/runners/<rid>/mode} & Simulazione o reale, a runner
  fermo; il reale chiede \cod{confirm}. \\
\rotta{POST}{/api/runners/<rid>/start} & Avvia. Con ordini reali chiede
  \cod{confirm\_real}; senza \cod{interval}, la cadenza viene dal
  timeframe. \\
\rigaTenue
\rotta{POST}{/api/runners/<rid>/stop} & Ferma il runner: gli ordini non
  ancora inviati del ciclo in corso non partono. \\
\rotta{GET}{/api/runners/<rid>/journal} & Le ultime righe del registro
  (\cod{limit}, di default \num{300}). \\
\rigaTenue
\rotta{POST}{/api/runners/<rid>/kill-switch/reset} & Riapre il conto
  chiuso dal kill-switch (il file del conto legato al runner), a runner
  fermo. \\
\rotta{POST}{/api/live/run\_cycle} & Un ciclo solo, per le prove da script.
  Su un broker vero sempre in simulazione. \\
\rigaTenue
\rotta{POST}{/api/live/start} & Ciclo continuo per le prove da script,
  con la stessa regola. \\
\rotta{GET}{/api/live/status} & Stato, numero di cicli, configurazione e le
  ultime \num{60} righe del ciclo di prova. \\
\bottomrule
\end{tabular}
\end{center}

\flushbottom
